\documentclass[tikz,border=10pt]{standalone}
\usepackage{tikz}
\usepackage{amsmath}
\usepackage{amssymb}
\usepackage{bm} 
\usetikzlibrary{positioning, arrows.meta, calc, decorations.pathreplacing}

\providecommand{\vect}[1]{\bm{#1}}

\begin{document}
\begin{tikzpicture}[
    >={Latex[length=5mm, width=3mm]}, % Frecce ulteriormente ingrandite per proporzione
    font=\huge, % Testo globale (linee/frecce) aumentato a \huge
    block/.style={
        rectangle,
        draw=black,
        thick,
        minimum width=4.5cm, 
        minimum height=2.2cm, 
        align=center,
        font=\Huge % Testo dei blocchi al massimo (\Huge)
    },
    bold_block/.style={
        rectangle,
        draw=black, 
        thick,
        font=\bfseries\Huge, % Testo dei blocchi al massimo e in grassetto
        minimum width=4.5cm,
        minimum height=2.2cm,
        align=center
    },
    image_block/.style={
        rectangle,
        draw=black,
        thick,
        minimum width=5cm,
        minimum height=3.5cm,
        align=center
    },
    arrow/.style={
        ->,
        thick,
        line width=1.5pt % Spessore linea aumentato
    }
]

% Nodes
\node[block] (observer) {Observer};
\node[bold_block, right=5.5cm of observer] (ppo) {PPO/\\SPPO}; % Spazio aumentato per ospitare la matrice grande
\node[block, right=5cm of ppo] (action) {Action};

% Drone image node
\node[image_block, right=6.5cm of action] (drone) {\includegraphics[width=4.5cm, height=3cm, keepaspectratio]{cad_model.png}}; 

% Initial input
\coordinate (start) at ($(observer.west) - (3.5cm,0)$);
\draw[arrow] (start) -- (observer.west) node[midway, above=0.4cm] {$\vect{p}_d, \vect{q}_d$};

% Arrow Observer -> PPO
\draw[arrow] (observer.east) -- (ppo.west) node[midway, above=0.3cm] { 
    $\renewcommand{\arraystretch}{1.2}
    \begin{bmatrix}
    \vect{e}_p^b \\
    \vect{v}^b \\
    \vect{q} \\
    \vect{\omega} \\
    \vect{q}_e
    \end{bmatrix} \in \mathbb{R}^{17}$
};

% Arrow PPO -> Action
\draw[arrow] (ppo.east) -- (action.west) node[midway, above=0.5cm] {$\vect{a}_t \in [-1, 1]$}; 

% Arrow Action -> Drone (Rimossa la radice quadrata)
\draw[arrow] (action.east) -- (drone.west) node[midway, above=0.5cm] {$\varpi_i \in [16, 100]$ Hz}; 

% Feedback loop (Drone -> Observer)
\coordinate (feedback_y) at ($(observer.south) - (0, 3cm)$); % Loop abbassato per evitare sovrapposizioni
\draw[arrow] (drone.south) -- (drone.south |- feedback_y) -- (observer.south |- feedback_y) -- (observer.south);

\end{tikzpicture}
\end{document}